\documentclass{article}
\usepackage{graphicx} % Required for inserting images
\usepackage{amsmath}
\usepackage[normalem]{ulem}
\usepackage{enumitem, pifont}
\usepackage{amssymb}
\usepackage[colorlinks=true, urlcolor=blue]{hyperref}
\usepackage{listings}
\usepackage{xcolor}


\newcommand{\cmark}{\ding{51}}
\newcommand{\xmark}{\ding{55}}

\newlist{todolist}{itemize}{1}
\setlist[todolist]{label=$\square$}

\newcommand{\done}{\rlap{$\square$}{\raisebox{2pt}{\large \cmark}}}
\newcommand{\wontfix}{\rlap{$\square$}{\large \xmark}}

\definecolor{codegreen}{rgb}{0,0.6,0}
\definecolor{codegray}{rgb}{0.5,0.5,0.5}
\definecolor{codepurple}{rgb}{0.58,0,0.82}
\definecolor{backcolour}{rgb}{0.95,0.95,0.92}

\lstdefinestyle{pythonstyle}{
    backgroundcolor=\color{backcolour},   
    commentstyle=\color{codegreen},
    keywordstyle=\color{magenta},
    numberstyle=\tiny\color{codegray},
    stringstyle=\color{codepurple},
    basicstyle=\ttfamily\footnotesize,
    breaklines=true,                 
    keepspaces=true,                 
    numbers=left,                    
    numbersep=5pt,                  
    showstringspaces=false,
    tabsize=4
}

\lstset{style=pythonstyle}

\title{2026 Formula 1 Season}


\begin{document}

\maketitle

\section{World Champions}
\subsubsection{2026 Title:}
Lando Norris won the WDC in the 2025 season.
\subsection{2027 Title}
Kimi Antonelli won the WDC in the 2026 season.

\textbf{At the Mexico Grand Prix, Kimi Antonelli was officially crowned as the Champion of this season}

\textit{As expected, Max Verstappen showed the fans during the Singapore GP that he still needs to be taken into account and proved it again with his P21-P2 climb last weekend}, even though he is out of contention for the first time since 5 years.

The Sao Paulo Grand Prix delivered a rather interesting Qualifying Session, \sout{with the worst cars of the grid} with some surprising cars in Q3 - the Haas and the Aston Martin.
\section{Calendar}
Races still left as of now:
\begin{itemize}
    \item Las Vegas
    \item Qatar
    \item Abu Dhabi
\end{itemize}
\section{Standings}
Even though both the WDC and WCC are decided, these are the current standings:

\subsection{Current WDC standings:}
\begin{enumerate}
    \item Kimi Antonelli
    \item Max Verstappen
    \item George Russell
    \item Lewis Hamilton
    \item Charles Leclerc
    \item Lando Norris
    \item Oscar Piastri
    \item Isack Hadjar
    \item Liam Lawson
    \item Pierre Gasly
\end{enumerate}

\subsection{Top 3 WCC standings}
The up to date WCC standings can be seen in Table~\ref{tab:data}.

\begin{table}[h]
       \centering

    \begin{tabular}{c c}
        \hline
        Team & Points \\
        \hline
        Mercedes & 652 \\
        Ferrari & 510 \\
        Red Bull Racing & 415 \\
        \hline
    \end{tabular}

    \caption{WCC standings.}
    \label{tab:data}
\end{table}

\section{Surprises}
Surprising events this season:
\begin{todolist}
    \item[\done] Ferrari mistake turned to upgrade
    \item[\done] Lando back-to-back win (Sao Paulo GP)
    \item[\done] George Russell Grand Slam
    \item[\wontfix] midfield team podium
    \item[\wontfix] Double DSQ
    \item[\done] 3 DNFs at home races
    \item[\done] Youngest WDC record beat
\end{todolist}

\section{Other}
Find more at: \href{http://colab.research.google.com/}{Google Colab}

If a driver finishes P5 in all races, his number of points will be:
\begin{lstlisting}[language=Python]
x = 10
y = 23 * x
print(y)
\end{lstlisting}

\section{Physics in Motorsport}
Frictional force:
\[
F_{\text{friction}} = \mu \cdot R_{\text{car}}
\]

Simple formula for speed:
\[
v = \frac{s}{t}
\]

Kinetic energy for constant velocity:
\[
E_{k} = \frac{1}{2} M_{\text{car}} v^{2}
\]

Centripetal force acting on the car driving on a circle/curve:
\[
F_{\text{Cen}} = \frac{M_{\text{car}} v^{2}}{r_{\text{corner}}}
\]

Maximum cornering force of the tires:
\[
F_{\text{max}} = \sqrt{\mu_{\text{track}} g M_{\text{car}}}
\]

Maximum cornering speed:
\[
v_{\text{max}} = \sqrt{\mu_{\text{track}} r_{\text{corner}} g}
\]
\includegraphics[width=\linewidth]{wykres.png}

\end{document}
